\documentclass[11pt,a4paper]{article}
\usepackage[T1]{fontenc}
\usepackage[utf8]{inputenc}
\usepackage{amsmath,amssymb,amsthm}
\usepackage[margin=1in]{geometry}
\usepackage[colorlinks=true,linkcolor=blue,urlcolor=blue]{hyperref}
\theoremstyle{plain}
\newtheorem{theorem}{Theorem}
\newtheorem{lemma}[theorem]{Lemma}
\newtheorem{proposition}[theorem]{Proposition}
\newtheorem{corollary}[theorem]{Corollary}
\theoremstyle{definition}
\newtheorem{remark}[theorem]{Remark}

\title{The reduction of OEIS A079144 modulo $k$:\\ a proof that the period divides
$\varphi(k)$}
\author{Adrian Perez Fontelles\\ \small Independent researcher}
\date{31 August 2026}
\begin{document}
\maketitle

\begin{abstract}
OEIS A079144 carries a conjecture of P.~Bala that for every positive integer $k$ the
sequence reduced modulo $k$ is eventually periodic with period dividing $\varphi(k)$.
It is true. The entry's own exponential generating function becomes, under the
substitution $t=e^{x}-1$, an ordinary power series with \emph{integer} coefficients;
that integrality forces all but finitely many terms of the Stirling expansion of $a(n)$
to vanish modulo $k$, and what is left is an integer combination of the functions
$n\mapsto i^{n}$ with $0\le i<k$. Each of those is eventually periodic with period
dividing $\varphi(k)$, and so is the sum. The argument also gives the exact point at
which periodicity starts: $n\ge v(k)$, where $v(k)$ is the largest exponent in the
prime factorisation of $k$.
\end{abstract}

\noindent\small 2020 Mathematics Subject Classification. 11B50, 11B73, 05A15.\normalsize

\section{The sequence and the conjecture}

OEIS A079144 is ``Number of labeled interval orders on n elements: (2+2)-free posets.''. It has offset $0$ and begins
\[
1, 1, 3, 19, 207, 3451, 81663, 2602699,\ \dots
\]
The entry states, not as a conjecture,
\begin{quote}\small
a(n) = Sum\_{k=0..n} k!*Stirling2(n,k)*A138265(k).
\end{quote}
that is, writing $A(x)=\sum_{n\ge 0}a(n)x^{n}/n!$,
\begin{equation}\label{eq:egf}
A(x)\;=\;\sum_{k\ge 0} k!\,S(n,k)\,b(k),\quad b=\text{A138265}.
\end{equation}
This is the only input taken from the entry, and it was checked against every one of the
$17$ terms of the entry's DATA section before use, not against a sample.

The entry separately carries the following comment:
\begin{quote}\small
\textbf{Conjecture} 1) Taking the sequence modulo an integer k gives an eventually periodic sequence with period dividing phi(k). For example, the sequence taken modulo 8 begins [1, 1, 3, 3, 7, 3, 7, 3, 7,  ...] and appears to have a pre-period of length 3 and a period of length 2 = (1/2)*phi(8). --- \emph{Peter Bala}, Dec 26 2021
\end{quote}
As of the ``Last modified'' line on the live entry (May 30 2026 16:40:09, revision
71) the statement is still recorded as a conjecture, and nothing on the entry
records it as settled either way.

Written out, the conjecture asserts: for every integer $k\ge 1$ there are $n_0$ and a
positive integer $P$ with $P\mid\varphi(k)$ such that
\begin{equation}\label{eq:conj}
a(n+P)\;\equiv\;a(n)\pmod{k}\qquad\text{for all } n\ge n_0 .
\end{equation}
It is proved below, with $P=\varphi(k)$ and $n_0=v(k)$, the largest exponent occurring in
the prime factorisation of $k$.

\section{The generating function in the variable $t=e^{x}-1$}

Set $t=e^{x}-1$. Then $e^{x}=1+t$, and substituting into \eqref{eq:egf},
\[
e^{x}-1=t,
\]
so that
\begin{equation}\label{eq:G}
A(x)\;=\;G\bigl(e^{x}-1\bigr),\qquad
G(t)\;=\;\sum_{k\ge 0} b(k)\,t^{k},\qquad b=\text{A138265}.
\end{equation}

\begin{lemma}\label{lem:int}
$G$ is a power series in $t$ with integer coefficients.
\end{lemma}

\begin{proof}
$G$ is a power series whose coefficients are the terms of the integer sequence A138265.
\end{proof}

Write $G(t)=\sum_{j\ge0}c_j t^{j}$, so $c_j\in\mathbb{Z}$; here $c_k=b(k)$, the terms of A138265, and the
first few are
\[
c_0,\dots,c_8 \;=\; 1, 1, 1, 2, 5, 16, 61, 271, 1372.
\]

\section{The Stirling expansion, and what survives modulo $k$}

Let $S(n,j)$ denote the Stirling numbers of the second kind.

\begin{lemma}\label{lem:stir}
$a(n)=\sum_{j=0}^{n} c_j\, j!\, S(n,j)$ for every $n\ge 0$. In particular $a(n)\in\mathbb{Z}$.
\end{lemma}

\begin{proof}
The standard expansion $\bigl(e^{x}-1\bigr)^{j}
=\sum_{n\ge j} j!\,S(n,j)\,x^{n}/n!$ holds because $j!\,S(n,j)$ counts the surjections
from an $n$-set onto a $j$-set. Substituting it into
$A(x)=G(e^{x}-1)=\sum_j c_j\bigl(e^{x}-1\bigr)^{j}$ and reading off the coefficient of
$x^{n}/n!$ gives the claim; the sum is finite because $S(n,j)=0$ for $j>n$. Integrality is
then immediate from Lemma~\ref{lem:int}.
\end{proof}

\begin{lemma}\label{lem:trunc}
For every $k\ge 1$ and every $n$,
$a(n)\equiv \sum_{j=0}^{k-1} c_j\, j!\, S(n,j) \pmod{k}$.
\end{lemma}

\begin{proof}
For $j\ge k$ the integer $j!$ is a multiple of $k$, since $k$ is one of the factors
$1,2,\dots,j$. Hence every term of Lemma~\ref{lem:stir} with $j\ge k$ vanishes modulo
$k$; this uses $c_j\in\mathbb{Z}$, which is Lemma~\ref{lem:int}.
\end{proof}

This is the step at which integrality is not a convenience but the whole point: were the
$c_j$ merely rational, $c_j\,j!$ need not be divisible by $k$ for large $j$ and the sum
would not truncate.

\begin{lemma}\label{lem:exp}
There are integers $d_0,\dots,d_{k-1}$, namely
$d_i=\sum_{j=i}^{k-1}(-1)^{j-i}\binom{j}{i}c_j$, with
\[
a(n)\;\equiv\;\sum_{i=0}^{k-1} d_i\, i^{n} \pmod{k}\qquad\text{for all } n\ge 0 .
\]
\end{lemma}

\begin{proof}
Counting surjections from an $n$-set onto a $j$-set by inclusion and exclusion on the
image gives $j!\,S(n,j)=\sum_{i=0}^{j}(-1)^{j-i}\binom{j}{i}i^{n}$. Substituting this
into Lemma~\ref{lem:trunc} and exchanging the two finite sums collects the coefficient of
$i^{n}$ into $d_i$ as stated. Each $d_i$ is an integer because each $c_j$ is.
\end{proof}

\section{The period}

For $k\ge 2$ let $v(k)$ be the largest exponent in the prime factorisation of $k$, so
$k\mid m^{v(k)}$ whenever $m$ is divisible by every prime factor of $k$, and
$v(k)\le\log_2 k$.

\begin{lemma}\label{lem:power}
Let $k\ge2$ and $0\le i<k$. Then $i^{\,n+\varphi(k)}\equiv i^{\,n}\pmod{k}$ for every
$n\ge v(k)$.
\end{lemma}

\begin{proof}
Write $k=k_1k_2$, where $k_1=\prod_{p\mid k,\ p\mid i}p^{v_p(k)}$ and $k_2=k/k_1$; then
$\gcd(k_1,k_2)=1$ and $\gcd(i,k_2)=1$.

Modulo $k_1$: for each prime $p\mid k_1$ we have $p\mid i$, so
$v_p\bigl(i^{n}\bigr)=n\,v_p(i)\ge n\ge v(k)\ge v_p(k)$. Hence $k_1\mid i^{n}$ for
$n\ge v(k)$, and likewise $k_1\mid i^{n+\varphi(k)}$; both sides are $\equiv 0$.

Modulo $k_2$: $i$ is a unit, so by Euler's theorem $i^{\varphi(k_2)}\equiv 1$. Since
$\gcd(k_1,k_2)=1$ we have $\varphi(k)=\varphi(k_1)\varphi(k_2)$, so
$\varphi(k_2)\mid\varphi(k)$ and therefore $i^{\varphi(k)}\equiv 1 \pmod{k_2}$, giving
$i^{\,n+\varphi(k)}\equiv i^{\,n}$.

The two congruences and the Chinese remainder theorem give the claim modulo $k=k_1k_2$.
(For $i=0$ note $0^{n}=0$ for $n\ge 1$ and $v(k)\ge1$.)
\end{proof}

\begin{theorem}\label{thm:main}
Let $a(n)$ be OEIS A079144. For every integer $k\ge1$,
\[
a\bigl(n+\varphi(k)\bigr)\;\equiv\;a(n)\pmod{k}\qquad\text{for all } n\ge v(k),
\]
where $v(k)$ is the largest exponent in the prime factorisation of $k$ (and $v(1)=0$).
In particular the sequence $\bigl(a(n)\bmod k\bigr)_{n\ge0}$ is eventually periodic with
period dividing $\varphi(k)$, with a pre-period of length at most $\log_2 k$. This is the
conjecture of Section~1.
\end{theorem}

\begin{proof}
The case $k=1$ is trivial. For $k\ge2$, Lemma~\ref{lem:exp} writes $a(n)$ modulo $k$ as
$\sum_{i<k}d_i i^{n}$ with integer $d_i$ independent of $n$, and Lemma~\ref{lem:power}
gives $i^{\,n+\varphi(k)}\equiv i^{\,n}$ for each $i<k$ once $n\ge v(k)$. Summing the
$k$ congruences with the weights $d_i$ gives
$a(n+\varphi(k))\equiv a(n)\pmod k$ for $n\ge v(k)$. The true period of the eventually
periodic sequence is then a divisor of $\varphi(k)$.
\end{proof}

\begin{remark}
The bound $n\ge v(k)$ cannot be lowered in general: for $k$ a prime power $p^{e}$ the
term $i=p$ contributes $p^{n}$, which is not divisible by $p^{e}$ until $n=e=v(k)$.
\end{remark}

\section{Verification}

Three checks, each independent of the algebra above.

First, the series $G$ of \eqref{eq:G} was expanded and $\sum_j c_j j! S(n,j)$ compared
against the entry's DATA at every one of the $17$ published indices; the values
agree exactly, and every $c_j$ came out an integer. A series that reproduced only some of
the published terms would not have been used.

Second, for each $k$ from $2$ to $40$ the sequence $a(n)\bmod k$ was computed exactly
--- not sampled --- from the recurrence $S(n,j)=j\,S(n-1,j)+S(n-1,j-1)$ carried out in
$\mathbb{Z}/k\mathbb{Z}$ on the truncated state
$\bigl(S(n,0),\dots,S(n,k-1)\bigr)$ of Lemma~\ref{lem:trunc}, and its exact eventual
period found by detecting the repetition of that state. The result is tabulated below.
In every case the period divides $\varphi(k)$, as Theorem~\ref{thm:main} requires.

\begin{center}\small
\begin{tabular}{rrr@{\qquad\qquad}rrr}
$k$ & $\varphi(k)$ & period & $k$ & $\varphi(k)$ & period \\ \hline
$2$ & $1$ & $1$ & $22$ & $10$ & $10$ \\
$3$ & $2$ & $2$ & $23$ & $22$ & $22$ \\
$4$ & $2$ & $1$ & $24$ & $8$ & $2$ \\
$5$ & $4$ & $4$ & $25$ & $20$ & $20$ \\
$6$ & $2$ & $2$ & $26$ & $12$ & $12$ \\
$7$ & $6$ & $6$ & $27$ & $18$ & $18$ \\
$8$ & $4$ & $2$ & $28$ & $12$ & $6$ \\
$9$ & $6$ & $6$ & $29$ & $28$ & $28$ \\
$10$ & $4$ & $4$ & $30$ & $8$ & $4$ \\
$11$ & $10$ & $10$ & $31$ & $30$ & $30$ \\
$12$ & $4$ & $2$ & $32$ & $16$ & $4$ \\
$13$ & $12$ & $12$ & $33$ & $20$ & $10$ \\
$14$ & $6$ & $6$ & $34$ & $16$ & $16$ \\
$15$ & $8$ & $4$ & $35$ & $24$ & $12$ \\
$16$ & $8$ & $2$ & $36$ & $12$ & $6$ \\
$17$ & $16$ & $16$ & $37$ & $36$ & $36$ \\
$18$ & $6$ & $6$ & $38$ & $18$ & $18$ \\
$19$ & $18$ & $18$ & $39$ & $24$ & $12$ \\
$20$ & $8$ & $4$ & $40$ & $16$ & $4$ \\
$21$ & $12$ & $6$ &  & & \\
\end{tabular}
\end{center}

Third, the congruence $a(n+\varphi(k))\equiv a(n)\pmod k$ of Theorem~\ref{thm:main} was
checked directly on the published terms for every $k$ in that range and every $n\ge v(k)$
for which both indices are published; it holds without exception.

\begin{thebibliography}{9}
\bibitem{oeis} The OEIS Foundation, \emph{The On-Line Encyclopedia of Integer
Sequences}, \url{https://oeis.org}, 2026. Sequence A079144.
\bibitem{comtet} L.~Comtet, \emph{Advanced Combinatorics}, Reidel, 1974.
\bibitem{fs} P.~Flajolet and R.~Sedgewick, \emph{Analytic Combinatorics}, Cambridge
University Press, 2009.
\bibitem{stanley} R.~P.~Stanley, \emph{Enumerative Combinatorics}, Volume 1, 2nd ed.,
Cambridge University Press, 2011.
\bibitem{hw} G.~H.~Hardy and E.~M.~Wright, \emph{An Introduction to the Theory of
Numbers}, 6th ed., Oxford University Press, 2008.
\end{thebibliography}

\end{document}
